\section{本讲主线：scaling laws 在真实模型里的落地}

上一讲建立了 data/model/compute scaling 的基本形式，本讲转向更困难的落地问题：真实模型并不只改变参数量，learning rate、batch、schedule、optimizer、normalization 与架构细节也会变化。核心任务是区分哪些规律可跨 scale 迁移，哪些结论只是某个 recipe 和实验区间的局部拟合。


\begin{knowledgebox}{讲义补充：源 Slide 1 如何接续 Lecture 9}
Lecture 9 讲 scaling laws 的基本形式和 Kaplan/Chinchilla/IsoFLOPS 等方法；Lecture 11 则进入“真实公开模型如何做 scaling”。重点不再是公式本身，而是公开论文里的 recipes：如何调 batch/LR，如何用 WSD 降低 sweep 成本，muP 是否真的稳定，optimizer scaling 如何验证。
\end{knowledgebox}


\singleslide{slides-images/slide-002.jpg}{本讲动机：Chinchilla 是否可落地、如何降低 scaling-fit 成本、哪些 parametrization 更稳定。}{2}
这页把抽象 scaling law 转成 recipe 审计：不仅预测 loss，还要预测 LR、batch、schedule 与 architecture 在目标规模是否仍有效。
\begin{knowledgebox}{讲义补充：源 Slide 2 的三个现实问题}
第一，Chinchilla 的 compute-optimal 方法在真实模型论文中是否还能工作。第二，完整从头训练网格非常贵，是否能用 WSD、早停或复用曲线节省计算。第三，架构、optimizer、learning rate、batch size 这些 hyperparameters 是否随 scale 改变，能否通过小模型可靠外推。
\end{knowledgebox}


\singleslide{slides-images/slide-003.jpg}{Scaling in practice：从 2022 经典结果转向 2024--2026 公开模型的真实 recipes。}{3}
Slide 3 强调研究对象已经变化：近期论文披露 WSD、muP、optimizer 与 MoE scaling，提供比单一 Chinchilla exponent 更丰富的证据。
\begin{knowledgebox}{讲义补充：源 Slide 3 的历史意义}
Lecture 9 主要围绕 Kaplan、Chinchilla 等早期 scaling 研究；这页提醒我们，2024-2026 的公开模型开始公布更多 scaling 细节。Scaling law 已经从研究论文变成训练报告里的工程流程：MiniCPM、DeepSeek、Qwen、Kimi、Hunyuan、LLaMA 3、MiniMax 都在展示不同风格的 scaling recipe。
\end{knowledgebox}

\begin{figure}[H]
\centering
\includegraphics[width=0.86\textwidth]{slide-004.jpg}
\caption{Slide 4：Initialization、optimizers、learning rate、batch 等 hyperparameters 都可能 scale-sensitive。}
\figsource{CS336 Spring 2026 Lecture 11 官方幻灯片。}
\end{figure}

\begin{importantbox}{术语消化：scale-sensitive hyperparameters}
\begin{tabular}{L{0.24\textwidth}L{0.34\textwidth}L{0.32\textwidth}}
\toprule
对象 & 为什么 scale-sensitive & 课程中的检验方式 \\
\midrule
Initialization & 宽度、深度改变会改变 activation 和 update magnitude。 & muP/SP 对比、activation/update 条件。 \\
Learning rate & 最优 LR 可能随参数量、batch、数据量变化。 & 小模型 grid search 或 muP 外推。 \\
Batch size & critical batch 随 loss 和数据规模变化。 & batch-loss 曲线、StepFun/DeepSeek/Qwen fits。 \\
Optimizer & 不同 optimizer 需要不同超参和 scaling rule。 & Muon、AdamC、exotic optimizers 的 scale study。 \\
\bottomrule
\end{tabular}
\end{importantbox}


\singleslide{slides-images/slide-005.jpg}{公开案例路线：先读 MiniCPM，再用 DeepSeek 检查规律能否跨团队复现。}{5}
选择这两个案例是为了比较相似问题的不同实验设计，而不是做模型排名。
\begin{knowledgebox}{讲义补充：源 Slide 5 的案例选择}
MiniCPM 是小模型高性能路线，公开了细致的 scaling computation 和 muP 使用；DeepSeek 是更大规模高性能开源模型，展示了不使用 muP 时如何直接估计 batch/LR 和做 IsoFLOP 分析。二者形成对照：一个强调稳定参数化，一个强调直接经验拟合。
\end{knowledgebox}

